\chapter{Who may sign}

The security of the set is the independence of the
signers. M-of-N is a story until the M cannot be
collected from one apartment, one employer, or one
compromised seed.

\section{The floor}

\begin{itemize}
  \item At least three signers. Never N-of-N: lose
  one key, brick the set.
  \item Hardware wallets. Browser-extension keys
  are not signers, they are incidents waiting.
  \item Around half the set as threshold, then
  raise it for upgrades and large treasuries.
  \item Dedicated address per Safe. Same device is
  fine if you use a different index and label it.
  Same seed still couples them: seed theft hits
  every index.
\end{itemize}

High-value sets want an external required signer:
security partner, advisor, someone who does not
share your Slack.

\section{Independence}

Spread people, orgs, devices, manufacturers,
geography. Mix roles: not five copies of the same
founder. Enough technical signers that someone in
the quorum can actually read calldata.

Do not store two signer keys on one laptop and
call it distribution.

\section{Proposer is not signer}

On Safe, authorize a delegated proposer that
cannot sign. Hash tools read the Safe API. The API
only has the transaction after it is proposed.
Without a proposer, the first signer is flying
blind.

All signers sign only. A non-signer executes.
Never ``sign and execute'' for the last signature.
Verification tools do not give you the hashes you
think they do for that combined action.

\section{Nested Safes}

A Safe as owner of another Safe. Sometimes an
entity wants 1-of-N inside their nest. The nested
set can then rotate its own owners without the
parent's consent.

Usually do not do this on a protocol set unless
you can explain why in writing.

\section{You can refuse}

If the set is God-mode, N-of-N, no hardware, no
delay on upgrades, or ``we all use the same
laptop,'' do not join. Being signer number six on
a theater quorum does not make you safer. It makes
you liable.
